% Generated by scripts/build_reports.py; edit reviewed Markdown inputs.
% Compile with XeLaTeX.
% Options for packages loaded elsewhere
\PassOptionsToPackage{unicode}{hyperref}
\PassOptionsToPackage{hyphens}{url}
\PassOptionsToPackage{dvipsnames,svgnames,x11names}{xcolor}
\documentclass[
  11pt,
]{article}
\usepackage{xcolor}
\usepackage[margin=1in]{geometry}
\usepackage{amsmath,amssymb}
\setcounter{secnumdepth}{-\maxdimen} % remove section numbering
\usepackage{iftex}
\ifPDFTeX
  \usepackage[T1]{fontenc}
  \usepackage[utf8]{inputenc}
  \usepackage{textcomp} % provide euro and other symbols
\else % if luatex or xetex
  \usepackage{unicode-math} % this also loads fontspec
  \defaultfontfeatures{Scale=MatchLowercase}
  \defaultfontfeatures[\rmfamily]{Ligatures=TeX,Scale=1}
\fi
\usepackage{lmodern}
\ifPDFTeX\else
  % xetex/luatex font selection
\fi
% Use upquote if available, for straight quotes in verbatim environments
\IfFileExists{upquote.sty}{\usepackage{upquote}}{}
\IfFileExists{microtype.sty}{% use microtype if available
  \usepackage[]{microtype}
  \UseMicrotypeSet[protrusion]{basicmath} % disable protrusion for tt fonts
}{}
\makeatletter
\@ifundefined{KOMAClassName}{% if non-KOMA class
  \IfFileExists{parskip.sty}{%
    \usepackage{parskip}
  }{% else
    \setlength{\parindent}{0pt}
    \setlength{\parskip}{6pt plus 2pt minus 1pt}}
}{% if KOMA class
  \KOMAoptions{parskip=half}}
\makeatother
\usepackage{longtable,booktabs,array}
\usepackage{caption}
\captionsetup[table]{skip=6pt}
\newcounter{none} % for unnumbered tables
\usepackage{calc} % for calculating minipage widths
% Correct order of tables after \paragraph or \subparagraph
\usepackage{etoolbox}
\makeatletter
\patchcmd\longtable{\par}{\if@noskipsec\mbox{}\fi\par}{}{}
\makeatother
% Allow footnotes in longtable head/foot
\IfFileExists{footnotehyper.sty}{\usepackage{footnotehyper}}{\usepackage{footnote}}
\makesavenoteenv{longtable}
\setlength{\emergencystretch}{3em} % prevent overfull lines
\providecommand{\tightlist}{%
  \setlength{\itemsep}{0pt}\setlength{\parskip}{0pt}}
\usepackage{amsmath,amssymb}
\usepackage{xurl}
\usepackage{fancyhdr}
\pagestyle{fancy}
\fancyhf{}
\fancyhead[L]{\small Pure Thought | Revised research specifications}
\fancyfoot[L]{\small 2026-09-17}
\fancyfoot[R]{\small\thepage}
\setlength{\headheight}{15pt}
\setlength{\emergencystretch}{3em}
\allowdisplaybreaks
\usepackage{bookmark}
\IfFileExists{xurl.sty}{\usepackage{xurl}}{} % add URL line breaks if available
\urlstyle{same}
% fallback for those not using the hyperref driver hyperxmp:
\makeatletter
\@ifundefined{xmpquote}{\newcommand{\xmpquote}[1]{#1}}{}
\makeatother
\hypersetup{
  pdftitle={S07. Finite-size criteria for parent-Hamiltonian spectral gaps},
  pdfauthor={Pure Thought research portfolio},
  colorlinks=true,
  linkcolor={blue},
  filecolor={Maroon},
  citecolor={Blue},
  urlcolor={blue},
  pdfcreator={LaTeX via pandoc}}

\title{S07. Finite-size criteria for parent-Hamiltonian spectral gaps}
\author{Pure Thought research portfolio}
\date{Revised 17 September 2026}

\begin{document}
\maketitle
\thispagestyle{empty}
\clearpage

\section{S07. Finite-size criteria for parent-Hamiltonian spectral gaps}

\textbf{Source:} synthesis challenge 23;
\href{https://github.com/rick-stevens-ai/Pure-Thought/blob/main/NOTES/30-Most-Compelling-Pure-Thought-AI-Challenges.md}{original
synthesis}. This topic is distinct from, or adds a distinct objective
to, the canonical PRD catalog.
\href{https://github.com/rick-stevens-ai/Pure-Thought/blob/main/reviews/SYNTHESIS-CROSSWALK.md}{Crosswalk}

\subsection{Plain-language guide}

\subsubsection{The problem in everyday terms}

A small quantum chain can have a clear energy separation above its
lowest state while that separation shrinks away as the chain grows. This
problem asks when an exact calculation on a short segment can rule out
that collapse for an entire specified family of longer chains. A theorem
connecting local and global behavior is the essential step.

\subsubsection{Key terms}

\begin{itemize}
\item
  \textbf{Quantum chain} --- A one-dimensional sequence of interacting
  quantum degrees of freedom.
\item
  \textbf{Matrix product state} --- A structured representation of a
  many-body state using products of small matrices.
\item
  \textbf{Parent Hamiltonian} --- A Hamiltonian constructed so that a
  specified state belongs to its ground-state space.
\item
  \textbf{Ground-state space} --- The collection of states attaining the
  minimum energy.
\item
  \textbf{Frustration-free} --- Having ground states that simultaneously
  minimize every local interaction term.
\item
  \textbf{Spectral gap} --- The energy separation above the entire
  ground-state space.
\item
  \textbf{Uniform gap} --- A positive lower bound on that gap that does
  not vanish as chain length increases.
\item
  \textbf{Thermodynamic limit} --- The limit of arbitrarily large
  systems.
\item
  \textbf{Finite-size criterion} --- A theorem translating a suitable
  local calculation into a conclusion for larger systems.
\end{itemize}

\subsubsection{Why this matters}

The gap helps distinguish low-energy behavior and enters rigorous
statements about correlations and stability under suitable conditions.
Proving a uniform gap turns a finite numerical observation into a
statement about a many-body family, which is often the scientifically
relevant object.

\subsubsection{What progress would mean}

Progress would verify all assumptions of a chosen finite-size theorem
and obtain a positive size-independent lower bound. Failing its
threshold would leave the question unresolved rather than prove
gaplessness. The result would apply to the specified family, not solve
the general spectral-gap problem.

\subsection{Revised core question}

For a specified frustration-free one-dimensional parent-Hamiltonian
family, can a validated local spectral calculation imply a uniform
many-body gap?

\subsection{Precise scope, target and outputs}

Fix exact MPS tensors, on-site dimension, interaction range, boundary
convention, and local positive projectors h\_i. Verify frustration
freeness and ground-space assumptions. Select a precise Knabe-type or
martingale theorem appropriate to that model. Compute interval
enclosures for the required finite-segment nonzero gaps and verify the
theorem's threshold to derive a uniform lower bound for the stated
sizes/boundaries. Degenerate ground spaces require identifying the first
strictly positive eigenvalue rather than assuming a unique ground state.

\subsection{Detailed critique of the synthesis framing}

The synthesis's jump from MPS to arbitrary PEPS makes a tractable
finite-size criterion look like a general spectral-gap decision method.
Such general algorithms are obstructed by undecidability results. A
positive finite-chain gap alone says nothing uniform as size grows; one
must meet a theorem's hypotheses and threshold. The revision targets
certificate automation for an explicit restricted family and preserves
unresolved outcomes when the chosen sufficient criterion fails.

\subsection{Validation and milestones}

Baseline: certify a standard parent model with a valid finite-size
theorem. Strong: a uniform parameter interval with an explicit gap lower
bound. Test a deliberately gapless/threshold-failing control without
claiming that criterion failure proves gaplessness.

Use the shared
\href{https://github.com/rick-stevens-ai/Pure-Thought/blob/main/reviews/EVIDENCE-STANDARD.md}{evidence
standard}. Unsuccessful search is unresolved unless a complete exclusion
proof is supplied. These are proposed baseline specifications, not
completed scientific results.

\href{https://arxiv.org/abs/1502.04573}{Cubitt, Pérez-García and Wolf,
Undecidability of the Spectral Gap} motivates the restriction to a
specified certifiable family.

\subsection{Displayed formulation and claim boundary}

\href{https://github.com/rick-stevens-ai/Pure-Thought/blob/main/supplemental/S07.md}{Supplemental
specification}.

\[H_N=\sum_i h_i,\qquad h_i\succeq0,\qquad
\Delta_N=\min(\operatorname{spec}H_N\setminus\{0\}),\]

\[\text{target: }\inf_{N\ \mathrm{in\ the\ stated\ family}}\Delta_N\ge\delta>0.\]

Verify frustration freeness and ground-space assumptions. A specified
finite-size gap theorem must supply the implication from validated local
gaps to the displayed uniform bound. Positive gaps on finitely many
chains are insufficient; failure of a sufficient criterion is not a
gaplessness proof.

\section{Evidence and execution standard}

This standard applies to the revised specifications. It is a proposed
contract for future research runs; this editorial revision does not
supply the scientific certificates themselves.

\subsection{What ``pure thought'' means here}

Use explicitly stated mathematical models and deductions without fitting
conclusions to experimental outcomes. Synthetic examples, exact model
inputs, analytic references, numerical search, and established software
are allowed. Physical constants, electronic-structure approximations,
constitutive laws, and empirical parameter sets must be identified when
used. A model-derived result is conditional on that model.

The original ban on legacy software conflicts with its proposed reliance
on NumPy, SciPy, CVXPY, nauty, and other libraries. Reimplementing every
optimizer and algebra system is not necessary for scientific
independence. Prefer a small independently checkable certificate and a
documented trusted base. A fresh implementation can be a project
objective, but must not be confused with a mathematical proof or novelty
claim.

\subsection{Four evidence levels}

{\def\LTcaptype{none} % do not increment counter
\begin{longtable}[]{@{}
  >{\raggedright\arraybackslash}p{(\linewidth - 4\tabcolsep) * \real{0.3333}}
  >{\raggedright\arraybackslash}p{(\linewidth - 4\tabcolsep) * \real{0.3333}}
  >{\raggedright\arraybackslash}p{(\linewidth - 4\tabcolsep) * \real{0.3333}}@{}}
\toprule\noalign{}
\begin{minipage}[b]{\linewidth}\raggedright
Label
\end{minipage} & \begin{minipage}[b]{\linewidth}\raggedright
What it establishes
\end{minipage} & \begin{minipage}[b]{\linewidth}\raggedright
What it does not establish
\end{minipage} \\
\midrule\noalign{}
\endhead
\bottomrule\noalign{}
\endlastfoot
Numerical evidence & A reproducible computation at stated precision,
sample size, and discretization & Exact truth, continuum coverage, or a
universal theorem \\
Finite-model certificate & An exact/validated statement about the
specified finite model or finite relaxation & An omitted tail,
thermodynamic limit, or physical realization \\
Conditional mathematical theorem & A proof over the stated domain with
every hypothesis and error controlled & Truth outside those assumptions
or a demonstrated physical system \\
Formal proof & A checked theorem in a named proof assistant, with
imported axioms/dependencies listed & Correct physical modeling merely
by virtue of formalization \\
\end{longtable}
}

Formalization is an optional extra deliverable unless the specific
project makes it mandatory. High precision, JSON output, solver status
and small residuals do not by themselves confer any of these proof
levels.

\subsection{Claim directions}

For minimization over a larger feasible set, the optimum is a lower
bound on the original minimum. A physical feasible point supplies an
upper bound. For maximization the directions reverse. Track feasibility
of the primal/dual witnesses separately from their objective values.

Infeasibility of an inner ansatz excludes only that ansatz.
Infeasibility of an outer relaxation can exclude the original problem if
the map into the relaxation is proved and all approximations preserve
necessity. Passing necessary tests means ``not excluded''; existence
needs a constructive sufficient argument.

A smallest observed example is not a minimality theorem. A complete
catalog must define the universe, equivalence relation, boundaries, and
proof that pruning misses nothing. A finite-size numerical trend does
not prove an asymptotic limit.

\subsection{Minimum run manifest}

Before calculation, record the problem ID and revision hash, exact
inputs and conventions, quantified domain, equivalence relation, target
observable, accepted evidence level, tolerances, and resource limits.
Record dependencies and precision/rounding mode. Where randomization is
used, preserve seeds, trial counts and raw sufficient statistics.
Resource limits are planning decisions, not evidence that a mathematical
solution cannot exist.

An output bundle should include:

\begin{itemize}
\tightlist
\item
  The statement being claimed, including all assumptions and
  quantifiers.
\item
  Construction or proof object, plus an independently runnable verifier
  and its trust assumptions.
\item
  Numerical/error ledger separating roundoff, discretization,
  truncation, optimization, sampling and modeling error where relevant.
\item
  Positive reference cases and negative controls that the verifier must
  reject.
\item
  Coverage: included domain, omitted domain, unresolved regions and
  termination reason.
\end{itemize}

These are requirements for future implementations, not a promise of
polynomial-size certificates for every problem.

\subsection{Checking specific certificate types}

\textbf{LP/SDP/SOS:} verify equality constraints, objective direction,
and cone feasibility using exact arithmetic or outward-rounded
intervals. Rationalizing a numerical solution can break PSD constraints.
Correct residuals with justified norm bounds; do not silently project
and reuse the old objective. Distinguish a proof of a polynomial
approximation from a proof for the original function.

\textbf{Interval numerics:} state function enclosures, domains and
outward rounding. A small root box needs existence and, if claimed,
uniqueness. Cover the complement for a global enumeration. Include
singularities and tails rather than discarding them numerically.

\textbf{SAT/SMT:} a solver's UNSAT label is not a certificate. Preserve
a proof supported by a checker and the proof that the encoding, symmetry
reduction and blocked solutions match the mathematical universe. DRAT
applies to an appropriate Boolean CNF proof, not automatically to every
real-arithmetic or SDP claim.

\textbf{Statistics:} specify the estimand, dependence assumptions,
confidence procedure and bias. Zero observed failures does not prove
zero failure probability. MCMC error bars based on estimated
autocorrelation are not generally rigorous mixing guarantees.

\subsection{Milestones and novelty}

Baseline means a working, independently validated result on a concrete
reference problem. Strong means a broader or sharper result with its
claimed evidence level. Research extension means a possible novel result
whose feasibility and significance need review. A negative theorem can
be valuable; an unsuccessful search can also be useful if its scope is
documented, but the two are not interchangeable.

Do not promise ``publication-quality'' in a fixed number of months.
Before claiming novelty, conduct a focused current literature comparison
against the exact theorem/model/metric. The sources in this revision
check key definitions and limitations; they are not an exhaustive
novelty audit across all 30 fields.

\end{document}
